% !TEX root = ../main.tex
\chapter{Introducción}
\label{chap:introduccion}

Cada mañana, en cada tienda de una cadena de distribución, alguien decide cuántas unidades de cada
producto fresco va a pedir. Lo hace sin conocer la demanda que tendrá que atender, y equivocarse
cuesta dinero en las dos direcciones: si pide de menos pierde la venta y el cliente se encuentra el
hueco en el lineal; si pide de más, el producto caduca y termina destruido. Esa decisión, repetida
miles de veces al día sobre artículos de vida útil corta y demanda volátil, es el problema que
aborda este trabajo. La gestión de la cadena de suministro en el sector \textit{retail} ha
evolucionado de un enfoque reactivo a uno proactivo apoyado en la analítica avanzada, pero la
previsión de la demanda de frescos sigue siendo uno de sus desafíos abiertos.

\section{Motivación}

\subsection{Motivación profesional}
La relevancia del problema no es solo operativa, sino también económica y ética. A nivel europeo, el
desperdicio alimentario asciende a unos 59 millones de toneladas anuales, con un coste estimado de
132.000 millones de euros para el conjunto de la cadena alimentaria
\cite{eucommission2023foodwaste}. En España, estas cifras han servido de base para una
transformación del marco legislativo con la Ley 1/2025 de Prevención de las Pérdidas y el
Desperdicio Alimentario \cite{leydesperdicio2025}, cuyas obligaciones para los agentes de la cadena
alimentaria rigen desde el 2 de abril de 2026. Entre ellas figura disponer de un plan de prevención
que concrete cómo se aplicará la jerarquía de usos de los excedentes, y su incumplimiento se
sanciona con multas de hasta 500.000 euros en las infracciones muy graves. Quedan excluidas las
microempresas y los establecimientos de menos de 1.300 metros cuadrados, de modo que la obligación
recae precisamente sobre las grandes cadenas de distribución.

En este contexto, si bien la legislación regula qué hacer con el producto sobrante, priorizando la
donación para consumo humano antes que la alimentación animal, el aprovechamiento industrial o el
compostaje, la gestión logística de esos excedentes sigue suponiendo un coste para la empresa. Por
tanto, la herramienta tecnológica más eficaz para cumplir con la ley y proteger al mismo tiempo el
margen de beneficio es la prevención en origen, que es el propósito exacto del Sistema de Apoyo a la
Decisión desarrollado en este proyecto: evitar el exceso de inventario antes de que llegue a
producirse.

Este Trabajo de Fin de Grado nace de la necesidad de alinear los modelos matemáticos de previsión
con la realidad de la toma de decisiones en la reposición. El caso de uso central no es una
predicción abstracta de demanda, sino la construcción de un \textbf{Sistema de Apoyo a la Decisión
(DSS)} que responda a una pregunta operativa concreta del responsable de aprovisionamiento:
\emph{\textquestiondown cuántas unidades debe pedir hoy cada SKU para cubrir la demanda de los
próximos días?} Para contestarla, el sistema no solo estima la demanda esperada y cuantifica
matemáticamente su incertidumbre, sino que la traduce en una cantidad de pedido que pondera de forma
explícita el coste por rotura de stock (\textit{stockout}) frente al coste por sobrestock y
desperdicio, con un umbral de decisión (el fractil crítico) común al catálogo y una
curva de cuantiles propia de cada producto, contribuyendo así al Objetivo de Desarrollo Sostenible (ODS) 12: Producción y Consumo
Responsables \cite{ods12un}.

\subsection{Motivación personal}
El origen de este trabajo es anterior al propio Trabajo de Fin de Grado. Surgió mientras preparaba mi
candidatura a unas prácticas en el área de tecnología de una cadena de distribución alimentaria, para
la que quise construir por
mi cuenta un prototipo de previsión de demanda que sirviese como muestra de trabajo. Aquel prototipo
nunca llegó a terminarse. Resolvía la parte cómoda del problema, entrenar un modelo y mirar su
error, pero se detenía justo antes de las preguntas incómodas, que son las que ocupan esta memoria.
El TFG ha sido la ocasión de retomarlo, completarlo y someterlo al rigor experimental del que
entonces carecía.

La orientación del trabajo responde a una convicción que se ha ido formando durante su desarrollo: un
modelo solo genera valor cuando alguien puede usarlo para decidir. No se trata de restar importancia
al fundamento teórico, que ocupa el núcleo metodológico de esta memoria, sino de exigirle que llegue
hasta una decisión ejecutable en lugar de agotarse en la evaluación. Un cuaderno de análisis que
produce una cifra correcta no es todavía una herramienta, y mi impresión es que esa distancia
entre el resultado experimental y el sistema en funcionamiento es precisamente la que más
interesa resolver en el entorno empresarial. Por eso el proyecto no se detiene en la tabla de
métricas y llega hasta la API, el cuadro de mandos y el despliegue, con la trazabilidad necesaria
para que cualquier número publicado pueda rehacerse.

Finalmente no me incorporé al sector de la distribución alimentaria, sino a una empresa de
automoción, donde he encontrado un problema de la misma familia planteado sobre otro catálogo. Allí
los vehículos se piden al fabricante con mucha antelación respecto a la venta y, cuando no encuentran
comprador, se traspasan al canal de ocasión como kilómetro cero. La estructura de la decisión es la
misma que se estudia en esta memoria: cuánto pedir hoy frente a una demanda que no se conocerá
hasta después, con un coste del exceso distinto del coste del defecto. Lo que cambia es la
naturaleza de esa asimetría. El excedente no se destruye, se degrada de canal, de modo que el coste
del sobrestock es una pérdida de margen y no una merma total, sobre un catálogo de pocas unidades de valor alto en
lugar de muchas de valor bajo. Esa coincidencia estructural ha reforzado el enfoque adoptado aquí: lo
que se transfiere entre sectores no es el modelo concreto ni la parametrización de sus costes, sino
la forma de encadenar la estimación de la demanda, la cuantificación de su incertidumbre y la
decisión de aprovisionamiento que de ambas se deriva.

\section{Definición del Problema}
El enfoque tradicional del \textit{demand forecasting} se ha centrado históricamente en la obtención de predicciones puntuales (\textit{point forecasts}) mediante la minimización de métricas de error promedio como el MAE (\textit{Mean Absolute Error}) o el RMSE (\textit{Root Mean Squared Error}). Aunque útiles, estos enfoques presentan severas limitaciones cuando intentan utilizarse aisladamente en entornos de \textit{retail} de producción:

\begin{enumerate}
    \item \textbf{Censura de datos por stockout:} Cuando un producto se agota en el estante, la demanda observada se detiene, pero la demanda latente (lo que los clientes habrían comprado) sigue existiendo. Musalem et al.~\cite{musalem2010structural} demuestran empíricamente que ignorar este fenómeno sesga a la baja la demanda estimada de los artículos afectados. El sesgo, además, es doble: como establecen Anupindi et al.~\cite{anupindi1998stockout}, las ventas de los productos en rotura quedan censuradas por la derecha, mientras que las de los productos disponibles quedan infladas por las sustituciones que absorben esa demanda desviada, de modo que ninguna de las dos series mide ya la demanda real. Si este efecto no se corrige, el sistema de reposición perpetuará un ciclo de ineficiencia.
    \item \textbf{Desconexión entre error estadístico y coste logístico:} En productos frescos, el coste de equivocarse ``por exceso'' (desperdicio) es estructuralmente distinto al coste de equivocarse ``por defecto'' (pérdida de venta). Ban y Rudin~\cite{ban2019big} demuestran en el contexto del problema del vendedor de periódicos (\textit{Newsvendor}) que la precisión predictiva y el coste de inventario pueden divergir: un modelo puede mejorar su error un 20\% sin mejorar el coste, o incluso empeorándolo, cuando la estructura de costes es marcadamente asimétrica. Su propuesta es resolver el problema en un solo paso, minimizando empíricamente el coste del \textit{Newsvendor}, en lugar de estimar primero la demanda y optimizar después. Este trabajo recorre deliberadamente la vía en dos etapas, y la pieza que la sostiene es la función de pérdida asimétrica o \textit{pinball loss}. Un modelo entrenado con el error cuadrático aprende a predecir la demanda \emph{media}; entrenado con la \textit{pinball loss} aprende a predecir un \emph{cuantil} concreto de la distribución de demanda, que es precisamente lo que la decisión de pedido necesita. Y existe una equivalencia analítica exacta entre el fractil crítico del \textit{Newsvendor} y el cuantil que esa pérdida minimiza, resultado que se remonta a la regresión cuantílica introducida por Koenker y Bassett~\cite{koenkerbassett1978} y sistematizada posteriormente por Koenker~\cite{koenker2005quantile}. Sobre la rejilla de cuantiles estimada, la capa de decisión localiza después el fractil que corresponde a cada producto. Esta separación tiene una ventaja operativa que el enfoque en un solo paso no ofrece: la estructura de costes puede cambiar, o ser distinta para cada producto, sin necesidad de reentrenar el modelo.
\end{enumerate}

Este trabajo aborda ambas limitaciones mediante un enfoque unificado: desde un \textbf{forecasting probabilístico} avanzado con garantías matemáticas (\textit{Mondrian Conformal Prediction}) hasta la traducción de esa distribución en una decisión de pedido mediante Investigación de Operaciones (modelo del \textit{Newsvendor}), todo ello empaquetado en una arquitectura MLOps lista para despliegue.

\section{Objetivos}
\label{sec:objetivos}
El objetivo general de este proyecto es diseñar, evaluar y desplegar un sistema completo de apoyo a la reposición para \textit{retail} fresco, elevando la solución desde un experimento estadístico hasta un sistema desplegable que cierra el bucle entre predicción y decisión de reposición, con el fin último de minimizar el desperdicio alimentario y maximizar el margen de beneficio operativo.

Para alcanzar este objetivo general, se plantean los siguientes objetivos específicos:
\begin{enumerate}[label=\textbf{O\arabic*.}, ref=O\arabic*]
    \item \label{obj:pipeline} Desarrollar un \textit{pipeline} de datos reproducible que incluya la imputación de demanda censurada a partir del dataset \textit{Fresh RetailNet}, verificando que ninguna observación censurada por stockout introduce fuga de información hacia la etiqueta objetivo.
    \item \label{obj:modelos} Implementar modelos globales de \textit{Machine Learning} probabilístico (CatBoost, LightGBM) capaces de compartir información entre productos, obteniendo un coste de inventario simulado en el conjunto de test inferior al de un \textit{baseline} estacional ingenuo (\textit{seasonal naïve}).
    \item \label{obj:conformal} Garantizar la calibración estadística de las predicciones mediante \textit{Mondrian Conformal Prediction} con optimización multi-objetivo (\textit{Pareto Frontier}), alcanzando una cobertura empírica $\geq 80\%$ para el intervalo $[q_{0{,}1}, q_{0{,}9}]$ (nivel $\alpha = 0{,}20$) sobre el conjunto de test.
    \item \label{obj:mlops} Operativizar el sistema mediante estándares \textit{MLOps} e Ingeniería de Software,
    incluyendo trazabilidad de experimentos en un almacén de corridas, explicabilidad (SHAP),
    una API operacional y un cuadro de mandos interactivo \textit{What-If} renderizado en servidor
    (Django), soportado por \textit{Docker} e Integración Continua (CI/CD).
\end{enumerate}

\section{Alcance}
El presente trabajo se enmarca en el contexto de la previsión de demanda a corto plazo para productos frescos en el sector \textit{retail}. A continuación se delimitan explícitamente los límites del sistema para evitar interpretaciones erróneas sobre su cobertura.

\textbf{Dentro del alcance:}
\begin{itemize}
    \item Previsión de la demanda acumulada a nivel SKU-tienda para un horizonte de planificación de corto plazo (días), alineado con el ciclo de reposición habitual en distribución de frescos.
    \item Datos de granularidad diaria provenientes del benchmark \textit{Fresh RetailNet}, que cubre 50.000 combinaciones tienda-producto a lo largo de 90 días de observación.
    \item Corrección del sesgo por demanda censurada (stockout) mediante técnicas de imputación supervisada.
    \item Cuantificación de la incertidumbre con garantías estadísticas formales (Mondrian Conformal Prediction) y optimización multi-objetivo de hiperparámetros.
    \item Traducción de la distribución de demanda predicha en una cantidad de pedido por producto y día, mediante el fractil crítico del \textit{Newsvendor}, y evaluación del coste logístico y del nivel de servicio que esa política induce.
    \item Operativización del sistema mediante estándares MLOps (trazabilidad de ejecuciones, API operacional, CI/CD) y visualización interactiva.
\end{itemize}

\textbf{Fuera del alcance:}
\begin{itemize}
    \item \textbf{Optimización de precios}: los precios y descuentos se tratan como variables exógenas observables, no como variables de decisión del sistema.
    \item \textbf{Cadena de suministro multi-escalón}: el sistema opera a nivel de tienda individual; la gestión de almacenes centrales, proveedores o logística de transporte queda excluida.
    \item \textbf{Predicción a medio y largo plazo}: no se aborda la planificación estacional ni la previsión más allá del horizonte de reposición inmediato.
    \item \textbf{Simulación de inventario con estado}: la evaluación económica resuelve una
    decisión independiente por origen de predicción. No se modelan el arrastre de existencias entre
    periodos, los pedidos en tránsito ni el \textit{lead time}, por lo que el sistema ordena
    políticas pero no reproduce un libro de reposición. Tampoco se aborda el reparto de una
    capacidad de almacén limitada entre productos cuando la suma de las cantidades individualmente
    óptimas la supera. Ambas extensiones se detallan en la Sección~\ref{sec:trabajo_futuro}.
    \item \textbf{Ejecución automática de pedidos}: el sistema genera recomendaciones de pedido cuantificadas, pero la decisión final y la emisión del pedido corresponden al responsable de aprovisionamiento.
    \item \textbf{Productos no perecederos}: la metodología está orientada a frescos con alta rotación y corta vida útil; su extrapolación a otras categorías requeriría validación adicional.
\end{itemize}

\section{Impacto Esperado}
El sistema no se dirige a un único perfil de usuario. Las decisiones de reposición implican a
personas con responsabilidades distintas, y el beneficio que cada una obtiene es también distinto:

\begin{itemize}
    \item \textbf{Responsable de aprovisionamiento} (usuario operativo principal). Recibe, para cada
    producto y cada día, una cantidad de pedido recomendada acompañada del intervalo de demanda que
    la justifica, en lugar de una cifra aislada sin indicación de su fiabilidad. Dispone además de un
    simulador que permite contrastar el efecto de modificar los supuestos de coste antes de
    comprometer el pedido, de modo que la herramienta apoya la decisión sin sustituirla.
    \item \textbf{Responsable de tienda u operaciones} (lectura agregada). Dispone del nivel de
    servicio y del coste agregados de la tienda, y de un listado de excepciones que señala los
    productos en riesgo, de modo que puede localizar dónde se concentran las roturas y las mermas
    sin revisar el catálogo entero. La agregación por categoría de producto queda como extensión
    natural: la categoría ya estructura la calibración conformal, pero el sistema no publica hoy
    coste ni nivel de servicio agrupados por ella.
    \item \textbf{Equipo técnico responsable del sistema}. Cada ejecución queda registrada con su
    configuración, sus métricas y la versión exacta del código que la produjo, y el sistema
    contempla el reentrenamiento periódico y la vigilancia de la degradación del modelo. El efecto
    esperado es reducir el coste de mantener el sistema en el tiempo, que es donde suelen fracasar
    los proyectos de analítica que sí funcionaron en su evaluación inicial.
    \item \textbf{Comunidad académica y docente}. La implementación completa, sobre un conjunto de
    datos público y con los experimentos reproducibles, sirve como referencia para un problema que
    la literatura suele abordar por partes separadas. El detalle de este legado se recoge en la
    Sección~\ref{sec:legado}.
\end{itemize}

Para la organización en su conjunto, el impacto se mide en las dos magnitudes que la operación
percibe: el coste logístico total y el nivel de servicio. Los resultados del
Capítulo~\ref{chap:resultados} cuantifican ambas mejoras y permiten estimar el efecto de trasladar
la metodología a un catálogo real. Ese mismo efecto tiene una lectura ambiental y regulatoria
directa, ya que el exceso de inventario en productos perecederos termina en merma: reducirlo en
origen es la vía más eficaz de contribuir al ODS~12 y de sustentar el Plan de Prevención que exige
la normativa vigente. El análisis detallado de la vinculación con los Objetivos de Desarrollo
Sostenible se recoge en el Anexo~\ref{anx:ods}.

\section{Aportaciones del Trabajo}
\label{sec:aportaciones}
Este trabajo no propone una técnica nueva, sino que establece cómo debe evaluarse la
recuperación de demanda censurada, y lo hace sobre un \textit{benchmark} publicado en mayo de
2025~\cite{freshretailnet2025}. Sus aportaciones son las siguientes.

\begin{enumerate}
    \item \textbf{El protocolo estándar para medir la reconstrucción de demanda censurada es
    circular, y ninguna métrica calculada sobre él ordena reglas de reconciliación.} El protocolo
    fabrica su verdad-terreno reduciendo la venta de días limpios en proporción a una rotura
    inventada, de modo que una regla que se limite a deshacer esa proporción recupera la
    respuesta exacta sin modelo alguno y obtiene el menor error de reconstrucción de todas las
    estrategias evaluadas, siendo la que menos garantía ofrece sobre una rotura real. El
    favoritismo alcanza igualmente a la evaluación por coste, que se puntúa sobre los mismos días.
    De ahí que la elección entre reglas se argumente aquí sobre bases de modelado, que la
    comparación de hiperparámetros dentro de una misma regla siga siendo válida, y que la ventaja
    medida para la corrección de censura se declare como cota superior
    (Sección~\ref{sec:imputacion_censura_sintetica}).

    \item \textbf{Una partición de validación puede decidir la respuesta de una búsqueda de
    hiperparámetros, y hacerlo en silencio.} El imputador reajusta sobre el panel que recibe, de
    modo que partir ese panel para obtener un conjunto de validación disjunto encoge también el
    conjunto de entrenamiento del maestro. Medida sobre un maestro así empobrecido, la
    sintonización parecía mejorar un 5{,}33\%; los mismos hiperparámetros resultaban un 12{,}44\%
    \emph{peores} que no sintonizar a la escala real de operación, de modo que el fichero
    persistido era activamente perjudicial allí donde se consumía. Separar por tiempo y aplicar la
    partición como una máscara sobre el panel completo, en lugar de como un recorte de él,
    invierte el veredicto: a escala de despliegue la sintonización mejora un 5{,}98\%, con el
    intervalo al 95\% íntegramente por debajo de cero y replicado sobre seis búsquedas
    independientes (Anexo~\ref{anx:tuning_imputador}).

    \item \textbf{En una evaluación por coste, el término de seguridad de la política decide el
    resultado.} Con un término común a todo el catálogo, conocer la demanda real sale más caro
    que la señal censurada sin corregir; con un término proporcional a cada serie, la métrica
    satura hasta no distinguir un imputador razonable de un oráculo. Solo sin término de
    seguridad el coste mide la señal, y la demanda real cuesta cero. El resultado es directamente
    reutilizable por cualquiera que compare señales de demanda por el coste que inducen
    (Sección~\ref{sec:resultados_gemelo}).

    \item \textbf{Medición de la calibración conformal estratificada sobre demanda
    reconstruida.} La variante \textit{Mondrian} se aplica aquí sobre un target reconstruido a
    partir de ventas censuradas, combinación que no cubren los trabajos que llevan la predicción
    conformal al \textit{Newsvendor} ni los que desarrollan la estratificación. Se mide qué le
    ocurre: la cobertura empírica queda por debajo del nivel nominal, la brecha se estrecha al
    ampliar el panel un orden de magnitud sin cerrarse, y el análisis atribuye el resto a la
    no-intercambiabilidad temporal entre calibración y test
    (Sección~\ref{sec:resultados_conformal}).

    \item \textbf{Sistema completo, trazable y reproducible.} Se implementa el recorrido entero
    desde el dato bruto hasta la decisión de pedido: preparación del panel, reconstrucción de la
    censura, ingeniería de características sin fuga temporal, modelos probabilísticos,
    calibración, política de inventario, API operacional, cuadro de mandos y despliegue
    contenerizado con integración continua. Cada ejecución queda registrada con su
    configuración, sus métricas y el \textit{commit} que la produjo, y toda cifra citada en esta
    memoria nombra la corrida que la generó. Esa trazabilidad es la que hace verificables las
    cuatro aportaciones anteriores.
\end{enumerate}

\section{Metodología}
El desarrollo del sistema sigue el ciclo de \textbf{CRISP-DM}~\cite{chapman2000crispdm}
(comprensión del problema, comprensión de los datos, preparación, modelado, evaluación y
despliegue) en su extensión para aprendizaje automático \textbf{CRISP-ML(Q)}~\cite{studer2021crispmlq},
cuya diferencia relevante aquí es que añade una fase final de monitorización y mantenimiento del
modelo desplegado, ausente del estándar original y necesaria porque el riesgo de degradación en un
entorno cambiante es precisamente uno de los fenómenos que este trabajo mide. A ello se suma una
particularidad que condiciona todas las etapas: al tratarse de datos temporales y de decisiones que se toman antes
de conocer la demanda, cada etapa debe respetar el orden del tiempo. Toda variable empleada en una
predicción ha de estar disponible en el instante en que esa predicción se emite, y toda comparación
entre alternativas debe realizarse sobre datos posteriores a los usados para ajustarlas. Este
principio, y no la elección de un algoritmo concreto, es el que gobierna el diseño experimental del
trabajo.

Sobre esa base, el trabajo se organiza en seis etapas encadenadas:

\begin{enumerate}
    \item \textbf{Análisis y preparación de los datos.} Caracterización del conjunto de partida,
    identificación de sus patrones y de sus patologías, y construcción de un panel diario homogéneo
    (Capítulo~\ref{chap:analisis_datos}).
    \item \textbf{Reconstrucción de la demanda censurada.} Estimación de la demanda que se habría
    producido en los días con rotura de stock, mediante tres estrategias alternativas que se
    comparan bajo un mismo protocolo en lugar de adoptarse por defecto
    (Capítulos~\ref{chap:metodologia} y~\ref{chap:resultados}).
    \item \textbf{Ingeniería de características y modelado probabilístico.} Codificación explícita
    de la historia y del calendario en variables sujetas a un contrato causal, y entrenamiento de
    modelos globales que predicen no un valor único sino varios cuantiles de la distribución de
    demanda (Capítulo~\ref{chap:metodologia}).
    \item \textbf{Calibración y validación temporal.} Ajuste de los intervalos de predicción para
    dotarlos de una garantía de cobertura verificable, y evaluación mediante \textit{backtesting}
    con ventanas móviles que reproducen el uso real del sistema
    (Capítulos~\ref{chap:metodologia} y~\ref{chap:resultados}).
    \item \textbf{Traducción de la predicción en decisión.} Conversión de la distribución estimada
    en una cantidad de pedido según la asimetría de costes del negocio, y evaluación del coste
    logístico y del nivel de servicio que esa política induce
    (Capítulos~\ref{chap:metodologia} y~\ref{chap:resultados}).
    \item \textbf{Operativización, monitorización y mantenimiento.} Empaquetado del sistema como
    aplicación desplegable con interfaz de consulta, cuadro de mandos y trazabilidad de cada
    ejecución, de modo que cualquier resultado publicado pueda reproducirse
    (Capítulo~\ref{chap:implementacion}). Esta etapa instancia la fase final de
    CRISP-ML(Q)~\cite{studer2021crispmlq}: detección de deriva sobre las variables de entrada,
    reentrenamiento periódico y una comparación de cadencias de reentreno que cuantifica si
    mantener el modelo al día compensa su coste.
\end{enumerate}

Las etapas no se recorrieron una sola vez ni en línea recta: varias decisiones de diseño se
revisaron al detectar que un resultado provenía de un defecto del protocolo de medida y no del
fenómeno estudiado. El criterio para cerrar cada etapa fue siempre el mismo: una diferencia entre
alternativas solo se considera real si se sostiene sobre datos no empleados para elegirlas y si su
magnitud supera la variabilidad del propio experimento, para lo cual toda comparación entre
alternativas se acompaña de un intervalo de confianza obtenido por remuestreo. El modelo final no
se elige por tanto con una métrica de error estadístico, sino leyendo conjuntamente las dos
magnitudes que la operación percibe, el coste logístico y el nivel de servicio, cada una con su
intervalo; el Capítulo~\ref{chap:resultados} expone qué ocurre cuando ambas no coinciden.

\section{Estructura del Documento}
La presente memoria se organiza en siete capítulos adicionales:
\begin{itemize}
    \item El \textbf{Capítulo 2: Estado del Arte} revisa la literatura actual sobre series temporales, modelos de \textit{boosting}, conformal prediction y gestión de inventarios.
    \item El \textbf{Capítulo 3: Análisis de Datos} describe el dataset utilizado y su preparación.
    \item El \textbf{Capítulo 4: Metodología} detalla la formulación matemática de los modelos, la calibración avanzada, la optimización multi-objetivo y la política de pedido.
    \item El \textbf{Capítulo 5: Implementación y Arquitectura} describe el diseño del software MLOps, la API de despliegue, el sistema de diagnóstico automático y la infraestructura CI/CD.
    \item El \textbf{Capítulo 6: Resultados} presenta el análisis comparativo de los modelos, la evaluación económica de la corrección de censura y la explicabilidad del sistema.
    \item El \textbf{Capítulo 7: Conclusiones} sintetiza los hallazgos principales.
    \item El \textbf{Capítulo 8: Presupuesto} detalla el coste económico del desarrollo del proyecto.
\end{itemize}

La memoria se completa con cuatro anexos. El \textbf{Anexo~\ref{anx:ods}} analiza la
vinculación del trabajo con los Objetivos de Desarrollo Sostenible de la Agenda 2030; el
\textbf{Anexo~\ref{anexo:dashboard}} recoge las capturas del cuadro de mandos web, que
documentan visualmente las vistas descritas en el Capítulo~\ref{chap:implementacion}; el
\textbf{Anexo~\ref{anx:tuning_imputador}} documenta la búsqueda de hiperparámetros del
imputador supervisado, cuyo resultado no modifica el sistema final pero deja una lección
metodológica sobre la validez de este tipo de búsquedas; y el \textbf{Anexo~\ref{anexo:mlflow}}
presenta la trazabilidad de corridas y artefactos en el almacén de experimentos de MLflow.

Por último, dado que el texto emplea terminología específica de la previsión probabilística y de la
gestión de inventarios, se ha incluido al inicio del documento, antes del índice, una
\textbf{Lista de Abreviaturas y Acrónimos} que reúne los términos y siglas utilizados a lo largo de
la memoria. Se recomienda su consulta al lector no familiarizado con el vocabulario del área.

\section{Convenciones}
Con el fin de facilitar la lectura, el texto sigue las siguientes convenciones tipográficas:

\begin{itemize}
    \item Los \textbf{términos técnicos procedentes del inglés} sin traducción asentada se marcan en
    cursiva al introducirse. Los de uso muy frecuente a lo largo de la memoria, como \textit{retail},
    \textit{forecasting} o \textit{stockout}, se emplean después en redonda para no
    recargar la lectura. La cursiva se reserva también para el énfasis puntual.
    \item Los \textbf{elementos de código} (nombres de módulos, ficheros, funciones, columnas de
    datos y claves de configuración) se escriben en tipografía monoespaciada, como en
    \texttt{observed\_demand} o \texttt{configs/experiment.yaml}. Esta tipografía no se emplea para
    ningún otro fin.
    \item Las \textbf{cifras} siguen la convención española: coma como separador decimal y punto como
    separador de millar. Los costes se expresan en \textbf{unidades monetarias} (u.m.) y no en euros,
    ya que el conjunto de datos utilizado no publica precios reales; las comparaciones entre
    políticas son por tanto relativas y no representan importes de un catálogo concreto.
    \item Los \textbf{acrónimos} se desarrollan la primera vez que aparecen, seguidos de la sigla
    entre paréntesis, y se recogen además en la Lista de Abreviaturas y Acrónimos del inicio del
    documento.
    \item Las \textbf{referencias bibliográficas} se citan de forma numerada entre corchetes y se
    detallan en la bibliografía final. Las direcciones web de herramientas, productos o repositorios
    se indican como nota al pie y no como entrada bibliográfica.
\end{itemize}
